% =============================================================
% บทที่ 3 - วิธีการดำเนินโครงงาน
% โครงงานประเภท B (งานวิจัย) — เน้น Data Collection, Preprocessing,
% Model, Evaluation ตามคู่มือ Guideline_Project_2026
% =============================================================

โครงงานนี้เป็นงานวิจัยเชิงเปรียบเทียบ จึงดำเนินงานตามลำดับของงานวิจัยด้านการค้นคืนสารสนเทศ คือเตรียมคลังข้อมูล สร้างดัชนีค้นหา พัฒนาสถาปัตยกรรมที่จะนำมาเปรียบเทียบ สร้างชุดทดสอบพร้อมเฉลย แล้วจึงวัดผลและทดสอบนัยสำคัญทางสถิติ

% -----------------------------------------------------------
\section{ขั้นตอนการดำเนินงาน}
% -----------------------------------------------------------

ทุกขั้นตอนทำด้วยโปรแกรมที่สั่งให้ทำงานซ้ำได้ตั้งแต่ต้นจนจบ ไม่มีขั้นตอนใดแก้ไขข้อมูลด้วยมือ ผู้อื่นจึงทำการทดลองซ้ำแล้วได้ผลเดิม ภาพรวมของเส้นทางข้อมูลทั้งหมดแสดงดังภาพที่ \ref{fig:pipeline}

% =============================================================
% ภาพเส้นทางข้อมูล — สองสายที่มาบรรจบกันที่การวัดผล
% วาดด้วย TikZ ไม่ต้องใช้โปรแกรมภายนอก
% =============================================================

\begin{figure}[H]
\centering
\begin{tikzpicture}[
  bx/.style  = {draw=black!55, rounded corners=2pt, minimum width=2.55cm,
                minimum height=1.25cm, align=center, inner sep=2pt, font=\scriptsize},
  dat/.style = {bx, fill=black!6},
  res/.style = {bx, fill=black!68, text=white},
  ev/.style  = {bx, fill=black!20, very thick, draw=black, minimum width=5.6cm,
                minimum height=1.0cm},
  lb/.style  = {font=\scriptsize\bfseries, anchor=west},
  ar/.style  = {-{Stealth[length=4pt]}, thick, black!65}
]

\def\dx{2.95}

% ---------- สายที่ 1 คลังสินค้า ----------
\node[lb] at (-1.45,1.12) {สายที่ 1 \quad คลังสินค้า};
\node[dat] (a1) at (0,0)        {ข้อมูลดิบ Amazon\\[1pt]\textbf{117{,}243} แถว\\16 คอลัมน์};
\node[dat] (a2) at (\dx,0)      {ล้างข้อมูล\\[1pt]\textbf{81{,}264} แถว\\12 คอลัมน์};
\node[dat] (a3) at (2*\dx,0)    {สกัดคุณลักษณะ\\[1pt]81{,}264 แถว\\\textbf{16} คอลัมน์};
\node[dat] (a4) at (3*\dx,0)    {ฝังเวกเตอร์\\[1pt]bge-m3\\\textbf{1{,}024} มิติ};
\node[res] (a5) at (4*\dx,0)    {ดัชนีค้นหา\\[1pt]Chroma\\ค้นได้สองภาษา};

\foreach \i/\j in {a1/a2, a2/a3, a3/a4, a4/a5} \draw[ar] (\i) -- (\j);

% ---------- สายที่ 2 ชุดทดสอบ ----------
\node[lb] at (-1.45,-1.08) {สายที่ 2 \quad ชุดทดสอบ};
\node[dat] (b1) at (0,-2.3)     {คำค้นจริงของผู้ใช้\\[1pt]ESCI + WANDS\\\textbf{1{,}980} คำ};
\node[dat] (b2) at (\dx,-2.3)   {ยิงเข้าคลัง\\[1pt]วัดคะแนน\\อันดับ 1};
\node[dat] (b3) at (2*\dx,-2.3) {คัดด้วยเกณฑ์ตัวเลข\\[1pt]\textbf{140} คำค้น\\3 กลุ่ม};
\node[dat] (b4) at (3*\dx,-2.3) {แปลเป็นภาษาไทย\\[1pt]คู่กันบรรทัด\\ต่อบรรทัด};
\node[res] (b5) at (4*\dx,-2.3) {ติดเฉลยด้วยมือ\\[1pt]เกณฑ์ ESCI\\\textbf{424} คู่};

\foreach \i/\j in {b1/b2, b2/b3, b3/b4, b4/b5} \draw[ar] (\i) -- (\j);

% สายที่ 1 ป้อนให้ขั้นยิงเข้าคลังของสายที่ 2
\draw[ar, dashed] (a5.south) -- ++(0,-0.52) -| (b2.north);

% ---------- บรรจบที่การวัดผล ----------
\node[ev] (ev) at (2*\dx,-4.3) {วัดผล 3 สถาปัตยกรรม $\times$ 2 ภาษา $=$ \textbf{840 คำตอบ}};
\draw[ar] (a5.south) -- ++(0,-3.46) -| (ev.east);
\draw[ar] (b5.south) |- (ev.east);
\draw[ar] (b3.south) -- ++(0,-1.0) -- (ev.north);

% ---------- เผยแพร่ ----------
\node[res, minimum width=3.3cm, minimum height=0.85cm] (hf) at (4*\dx,-4.3)
      {เผยแพร่บน HuggingFace};
\draw[ar] (b5.south) -- (hf.north);

\end{tikzpicture}
\caption{เส้นทางข้อมูลสองสายที่มาบรรจบกันที่ขั้นวัดผล}
\label{fig:pipeline}
\end{figure}


\subsection{การรวบรวมข้อมูล (Data Collection)}

ใช้ข้อมูลสินค้าจริงจากคลัง Amazon Reviews 2023 \cite{hou2024amazon} ผ่านชุดข้อมูลที่ผ่านการจัดรูปแบบแล้วบน HuggingFace ได้ข้อมูลตั้งต้น 117{,}243 แถว 16 คอลัมน์ ครอบคลุมสินค้าหลายหมวด เหตุผลที่ใช้ข้อมูลจริงแทนการสร้างข้อมูลจำลองคือชื่อสินค้าจริงมีลักษณะที่ข้อมูลจำลองไม่มี เช่น ชื่อยาว ตัวย่อทางการค้า และรหัสรุ่นสินค้า ซึ่งเป็นสิ่งที่ระบบค้นคืนต้องรับมือจริง

\subsection{การเตรียมข้อมูล (Data Preprocessing)}

ขั้นล้างข้อมูลใช้ข้อมูลครบทุกหมวดโดยไม่ตัดหมวดใดทิ้ง ดึงคอลัมน์หมวดสินค้าจากชื่อไฟล์ต้นทางเนื่องจากคอลัมน์เดิมขาดหายทั้งหมด และตัดแถวที่ไม่มีราคาออกเพราะโครงงานนี้ต้องประเมินความสามารถในการทำตามงบประมาณ เหลือข้อมูล 81{,}264 แถว 12 คอลัมน์

\subsection{การสกัดคุณลักษณะ (Feature Extraction)}

ขั้นนี้สร้างคอลัมน์ข้อมูลเพิ่มจากข้อมูลที่มีอยู่ ได้แก่ คะแนนความนิยมแบบเบย์ส คะแนนความคุ้มค่า ข้อความสำหรับฝังเวกเตอร์ และหมวดย่อย รวมเป็น 16 คอลัมน์ คุณลักษณะที่เป็นคุณสมบัติของสินค้า เช่น ยี่ห้อ สี และวัสดุ ไม่ได้แยกเป็นคอลัมน์ แต่รวมไว้ในข้อความที่ใช้ฝังเวกเตอร์ เพื่อให้แบบจำลองภาษาเป็นผู้ตีความแทนการเขียนกฎดึงคำ

\subsection{การฝังเวกเตอร์และสร้างดัชนีค้นหา}

แปลงข้อความสินค้าเป็นเวกเตอร์ด้วยแบบจำลอง \texttt{bge-m3} ขนาด 1{,}024 มิติ ซึ่งรองรับทั้งภาษาไทยและภาษาอังกฤษในแบบจำลองเดียว จากนั้นสร้างดัชนีค้นหาด้วยฐานข้อมูลเวกเตอร์ Chroma โดยเก็บข้อมูลประกอบที่ทุกสถาปัตยกรรมเรียกใช้ร่วมกัน

\subsection{การพัฒนาสถาปัตยกรรมที่นำมาเปรียบเทียบ}

พัฒนาโมดูลการทำงานเป็นชิ้นส่วนที่ใช้ร่วมกันทุกสถาปัตยกรรม ได้แก่ โมดูลความจำ โมดูลวางแผน โมดูลแปลภาษา โมดูลค้นคืน โมดูลจัดอันดับ โมดูลเขียนคำค้นใหม่ และโมดูลสร้างคำตอบ แล้วจึงประกอบเป็นสถาปัตยกรรม 3 แบบที่ต่างกันเฉพาะชั้นที่แบบจำลองภาษาเป็นผู้ตัดสินใจ การออกแบบเช่นนี้ทำให้การเปรียบเทียบยุติธรรม เพราะทุกแบบเรียกใช้โมดูลค้นคืนและโมดูลจัดอันดับตัวเดียวกัน

\subsection{การสร้างชุดทดสอบและเฉลย}

ยืมข้อความคำค้นจริงของผู้ใช้จากชุดข้อมูลมาตรฐานที่ตีพิมพ์แล้ว 2 ชุด นำมายิงเข้าคลังของโครงงานนี้จำนวน 1{,}980 คำ แล้วคัดด้วยเกณฑ์ตัวเลขเหลือ 140 คำค้น แบ่งเป็น 3 กลุ่มตามสิ่งที่คลังมี จากนั้นติดเฉลยความตรงของสินค้าด้วยมือตามเกณฑ์มาตรฐาน 4 ระดับ ชุดคำค้นถูกล็อกก่อนรันและห้ามแก้หลังเห็นผล

\subsection{การวัดผลและทดสอบนัยสำคัญ}

รันทั้ง 3 สถาปัตยกรรมบนชุดทดสอบเดียวกันทั้งสองภาษา วัดด้วยตัวชี้วัดที่ใช้เฉลยและตัวชี้วัดที่โปรแกรมตรวจเองได้ จากนั้นทดสอบว่าส่วนต่างที่วัดได้แยกจากความบังเอิญได้หรือไม่ ด้วยการทดสอบจับคู่แบบสลับเครื่องหมายที่นับครบทุกกรณี

% -----------------------------------------------------------
\section{เครื่องมือที่ใช้ในการดำเนินงาน}
% -----------------------------------------------------------

โครงงานนี้ประมวลผลบนเครื่องส่วนบุคคลทั้งหมด ไม่ได้เรียกใช้บริการแบบจำลองภาษาผ่านเครือข่าย เพื่อให้ควบคุมสภาพแวดล้อมการทดลองได้และทำซ้ำได้โดยไม่ขึ้นกับผู้ให้บริการภายนอก

\subsection{ฮาร์ดแวร์}

เครื่องมือด้านฮาร์ดแวร์ที่ใช้แสดงดังตารางที่ \ref{tab:hardware}

\begin{table}[H]
\centering
\caption{รายการฮาร์ดแวร์ที่ใช้ในการดำเนินงาน}
\label{tab:hardware}
\begin{tabular}{|c|l|l|}
\hline
\textbf{ลำดับ} & \textbf{รายการ} & \textbf{รายละเอียด} \\
\hline
1 & หน่วยประมวลผลกลาง & Intel Core i5-12400F (6 แกน 12 เธรด) \\ \hline
2 & หน่วยความจำหลัก & 16 GB \\ \hline
3 & หน่วยประมวลผลกราฟิก & NVIDIA GeForce RTX 4060 หน่วยความจำ 8 GB \\ \hline
4 & หน่วยเก็บข้อมูล & SSD NVMe ความจุ 500 GB \\
\hline
\end{tabular}
\end{table}

หน่วยความจำของหน่วยประมวลผลกราฟิกขนาด 8 GB เป็นข้อจำกัดที่กำหนดขนาดของแบบจำลองภาษาที่ใช้ได้ จึงเลือกแบบจำลองขนาด 8 พันล้านพารามิเตอร์ ซึ่งเป็นขนาดใหญ่ที่สุดที่ทำงานได้บนเครื่องนี้โดยไม่ต้องสลับหน่วยความจำ ข้อจำกัดนี้มีผลโดยตรงต่อข้อค้นพบของโครงงาน

\subsection{ซอฟต์แวร์}

เครื่องมือด้านซอฟต์แวร์ที่ใช้แสดงดังตารางที่ \ref{tab:software}

\begin{table}[H]
\centering
\caption{รายการซอฟต์แวร์ที่ใช้ในการดำเนินงาน}
\label{tab:software}
\small
\begin{tabular}{|c|l|c|p{5.2cm}|}
\hline
\textbf{ลำดับ} & \textbf{ซอฟต์แวร์} & \textbf{เวอร์ชัน} & \textbf{วัตถุประสงค์} \\
\hline
1 & Windows 11 Pro & build 26200 & ระบบปฏิบัติการ \\ \hline
2 & Python & 3.12.10 & ภาษาที่ใช้พัฒนาทั้งโครงงาน \\ \hline
3 & Ollama & 0.32.1 & รันแบบจำลองภาษาในเครื่อง \\ \hline
4 & Typhoon 2 (Llama 3.1) & 8B instruct & แบบจำลองภาษาสำหรับวางแผน แปล และเขียนคำตอบ \\ \hline
5 & BGE-M3 & -- & แบบจำลองฝังเวกเตอร์สองภาษา 1{,}024 มิติ \\ \hline
6 & Chroma & 1.5.9 & ฐานข้อมูลเวกเตอร์และดัชนีค้นหา \\ \hline
7 & pandas & 3.0.6 & จัดการข้อมูลแบบตาราง \\ \hline
8 & NumPy & 2.5.3 & คำนวณเชิงตัวเลข \\ \hline
9 & huggingface\_hub & 1.33.0 & ดึงข้อมูลต้นทางและเผยแพร่ชุดทดสอบ \\ \hline
10 & Visual Studio Code & -- & เครื่องมือเขียนโปรแกรม \\
\hline
\end{tabular}
\end{table}

\subsection{ขนาดข้อมูลที่ใช้}

ขนาดของข้อมูลในแต่ละขั้นแสดงดังตารางที่ \ref{tab:datasize}

\begin{table}[H]
\centering
\caption{ขนาดข้อมูลในแต่ละขั้นของการเตรียมข้อมูล}
\label{tab:datasize}
\begin{tabular}{|c|l|r|l|}
\hline
\textbf{ลำดับ} & \textbf{ขั้นตอน} & \textbf{จำนวน} & \textbf{หน่วย} \\
\hline
1 & ข้อมูลดิบ & 117{,}243 & แถว (16 คอลัมน์) \\ \hline
2 & หลังล้างข้อมูล & 81{,}264 & แถว (12 คอลัมน์) \\ \hline
3 & หลังสกัดคุณลักษณะ & 81{,}264 & แถว (16 คอลัมน์) \\ \hline
4 & เวกเตอร์ที่สร้าง & 81{,}264 & ชุด (1{,}024 มิติ) \\ \hline
5 & คำค้นจริงที่นำมาคัด & 1{,}980 & คำค้น \\ \hline
6 & ชุดทดสอบที่ล็อกแล้ว & 140 & คำค้น (ภาษาละ 140) \\ \hline
7 & เฉลยความตรงที่ใช้คิดผล & 424 & คู่ (คำค้น, สินค้า) \\
\hline
\end{tabular}
\end{table}

% -----------------------------------------------------------
\section{แผนการดำเนินงาน}
% -----------------------------------------------------------

แผนการดำเนินงานตลอดระยะเวลาโครงงานแสดงดังตารางที่ \ref{tab:gantt}

\begin{table}[H]
\centering
\caption{แผนการดำเนินงาน (Gantt Chart)}
\label{tab:gantt}
\small
\begin{tabular}{|c|p{3.9cm}|c|c|c|c|c|c|}
\hline
\textbf{ลำดับ} & \textbf{กิจกรรม} & \textbf{ด.1} & \textbf{ด.2} & \textbf{ด.3} & \textbf{ด.4} & \textbf{ด.5} & \textbf{ด.6} \\ \hline
1 & ทบทวนวรรณกรรม & $\bullet$ & $\bullet$ & & & & \\ \hline
2 & รวบรวมและเตรียมข้อมูล & $\bullet$ & $\bullet$ & & & & \\ \hline
3 & สร้างดัชนีค้นหา & & $\bullet$ & & & & \\ \hline
4 & พัฒนาโมดูลร่วมและ 3 สถาปัตยกรรม & & & $\bullet$ & $\bullet$ & & \\ \hline
5 & สร้างชุดทดสอบและติดเฉลย & & & & $\bullet$ & $\bullet$ & \\ \hline
6 & รันการทดลองและวัดผล & & & & & $\bullet$ & \\ \hline
7 & วิเคราะห์ผลและทดสอบนัยสำคัญ & & & & & $\bullet$ & $\bullet$ \\ \hline
8 & จัดทำรายงานและเผยแพร่ชุดข้อมูล & & & & & & $\bullet$ \\
\hline
\end{tabular}
\end{table}
